\section{A research programme toward quasi-polynomial recognition}
\label{sec:research}

The new parameter moves the main question from discovering all local
quadrilateral choices to controlling the places where geometric and
algebraic intersections disagree. This is useful because a cohomology
generator is inexpensive to compute, and the disagreement is a directly
testable geometric object. The next results should establish control of
that object under certified changes of representation.

\subsection{The structural theorem that would complete this route}

The most direct objective is the following concrete question.

\begin{quote}
Is there an explicit procedure, taking $n^{O(\log n)}$ bit time, which
produces certified polynomial-size exteriors of an $n$-crossing diagram,
such that an unknot has a standard vertex meridian with
$p+|F(D)|\le C\log(n+2)$ in at least one produced source, for an absolute
known constant $C$?
\end{quote}

An affirmative answer combines with Corollary~\ref{cor:qp}; nothing else
in that corollary requires expanding a normal surface. Several aspects
of the wording are essential. The favourable source must be found within
the bound, its relation to the original diagram must be certified, and
vertex potentials count in $p$. A construction with many vertices cannot
silently invoke the one-vertex cancellation theorem. A promising normal
disc that is not a standard vertex disc does not yet meet the promise.

This question is deliberately not elevated to a conjecture that all
sources satisfy. The Whitehead construction rules out that broader
statement. Even the restriction to all one-vertex triangulations remains
an open problem in this investigation and may be false. Constructively
choosing a family of representations is a more flexible target.

\subsection{Geometric questions, in order of dependence}

\begin{enumerate}[label=G\arabic*.]
\item \textbf{Can low cancellation be transferred to a vertex disc?}
Suppose some essential normal disc has cancellation on at most $k$
edges. Is there a standard vertex meridian with at most $f(k)$ such
edges, ideally $f(k)=O(k)$, after a controlled source modification?
The ordinary vertex-disc theorem does not include this support control.
Intersecting a sector with $W_F$ produces extreme rays of the intersection,
which need not be extreme in the full sector. This is exactly the gap a
proof must bridge. A useful first result would identify sufficient
conditions on oriented normal sums under which no cancellation is created
on an edge that was previously coherent.

\item \textbf{Can the Whitehead obstruction persist with one vertex?}
Construct a one-vertex family with $\kappaV=\Omega(t)$, or prove a
quantitative obstruction to collapsing the many interior loop vertices
in Theorem~\ref{thm:whitehead} without destroying the independent
geometric-index barriers. Either result would substantially narrow the
correct global claim. The geometric index of an individual embedded loop
does not by itself control a whole collapsed graph, whose edges may share
vertices and change their isotopy classes.

\item \textbf{How does cancellation behave under certified moves?}
For $2$--$3$, $3$--$2$, boundary layering and crushing operations, derive
local inequalities for the best achievable cancellation support of an
essential disc, while also controlling standard extremality. Boundary
layering has the exact favourable theorem proved here. General Pachner
moves can change the coherent baseline, so comparing only geometric
weights before and after a move is insufficient. A useful move record
should transport the cocycle, candidate disc, affected edge stars and
source correspondence together.

\item \textbf{Can repeated obstacles be removed in batches?}
The Whitehead example has many parallel copies of a single geometric
gadget. Develop a compressed quotient of repeated edge-star regions,
with binary multiplicities and a proof that essential-disc discovery is
preserved. A parameter counting independent gadget types might be small
where literal cancellation support is linear. The hard part is retaining
the ordered attachments and their fibre data; merely aggregating edge
weights can change the embedding problem. The interface operator of
Section~\ref{sec:parallelity-interface} supplies part of the data needed
for a safe batch transformation.

\item \textbf{Can a complexity potential force logarithmic depth?}
Seek a certified simplification or hierarchy move that either exposes a
meridian or decreases a nonnegative structural potential by a fixed
factor. The potential should account for core size, boundary pattern,
vertex-gauge dimension and independent cancellation regions. A decrease
in coordinate weight alone is unsuitable: the layered family has huge
weights but a constant discovery parameter. A decrease in tetrahedron
count alone may fail during useful Pachner exploration. A potential
theorem must bound branching as well as depth.
\end{enumerate}

\subsection{Algebraic and algorithmic questions}

\begin{enumerate}[label=A\arabic*.]
\item \textbf{Implement the general sparse-incidence theorem.}
The highest-confidence next coding task is the row-first enumeration in
Theorem~\ref{thm:incidence}, including its exceptional normalization case.
Use partial injective row-to-edge matchings, incremental exact rank and
canonical descriptions to avoid permutations of the same system.
Completeness should first be compared with the elementary algorithm at
$k=2,3$ on small sources. No heuristic pruning should be credited in the
asymptotic bound without a certificate that it preserves every relevant
invertible system.

\item \textbf{Can the exponent depend on interaction rather than support?}
The determining equations form a hypergraph with row size at most four.
Study dynamic programming over certified decompositions of the affected
edge-star interaction graph. A bag must retain compatible four-point
sectors, rational boundary weight relations and enough global information
to verify a primitive essential disc. Storing only possible total Euler
characteristics would be unsound. The goal is an explicit bound in both
the interaction width and the coefficient bit length, followed by a
geometric theorem controlling that width on selected sources.

\item \textbf{Can algebraic signatures reduce duplicate active systems?}
The present global list has many repeated hyperplanes, and one ray may
have many active descriptions. Investigate reverse search on active
bases or rank certificates indexed by sparse minors. A useful target is
polynomial delay per distinct candidate on a fixed $W_F$. The number of
distinct candidates can still be exponential in the dimension, so such
an output-sensitive result would improve implementation without alone
settling the global complexity question.

\item \textbf{Can the vertex gate replace more expensive positive checks?}
The exact minor and boundary parity certificate is complete for standard
vertex meridians. Measure its crossover with the existing exposed-unit
ray gate and compressed component observer. A sensible schedule first
uses cheap sufficient rank certificates, then the complete supported
rank test when a theoretical vertex-only mode is requested, and reserves
the richer observer for additional restricted rays. Source preparation
and reuse of matrix factors must be included in comparisons.

\item \textbf{How can multiple-vertex sign chambers be compressed?}
The current extension pays $t^{O(p)}$ for the arrangement of cocycle
hyperplanes. Examine whether a certified spanning-tree gauge, boundary
collapses or a small set of relevant vertex potential differences can
reduce $p$ without excluding the unknown disc. Fixing arbitrary vertex
potentials is not harmless: signed intersections of the actual disc may
use a different representative. A valid reduction must transport the
geometric problem, rather than simply delete gauge variables algebraically.
\end{enumerate}

\subsection{Compressed geometry and hierarchy questions}

\begin{enumerate}[label=H\arabic*.]
\item \textbf{Construct the boundary homeomorphisms, not only handle plans.}
For each full or partial disc-base attachment, output the actual collar
maps and transported marked arcs. Prove that disjoint source interfaces
remain disjoint and that the output records their induced order. The
correct relative handle count is already known here; an executable
topological replacement requires more data than that count.

\item \textbf{Convert the result to admissible handles.}
Show that the transported relative additions can be changed to the
uniform-type handle structures required by the hierarchy theory, with
polynomial encoded size. State precisely which boundary patterns are
preserved and which tests establish admissibility. A bounded inherited
core plus interval bundles does not automatically provide this step.

\item \textbf{Extend fibre information beyond separate boundary circles.}
The present signed double cover records transport around each marked
boundary circle. Non-disc bases may require paths connecting distinct
circles and coherent horizontal boundary markings. Identify a finite
set of such paths, bound its size, and certify its relation to the source
charts without expanding the base. Annular and M\"obius examples should
remain explicit regression cases.

\item \textbf{Close the orbit-engine complexity gap.}
The interface theorem has a polynomial reduction to ordinary AHT
counting; the implementation uses weighted traces. Either prove that
these concrete traces and their arithmetic stay polynomial in encoded
source size, or substitute a rigorously bounded counting implementation
for the proof-critical mode. Fast behaviour on long binary multiplicities
is valuable evidence, but does not control every possible interval system.
\end{enumerate}

\subsection{An experimental programme with useful failure criteria}

The next corpus should combine several independently generated classes:
small enumerated solid tori, diagram exteriors with verified provenance,
Pachner walks, layered sources, and sources containing controlled
geometric-index gadgets. For each, record $t$, vertex count, cocycle rank,
the best certified vertex cancellation support found, sector nullity,
maximum coordinate bit length and edge-star degree distribution.
Keep a distinct field for a lower bound or exhaustive result; a heuristic
failure to find a disc is not a lower bound on $\kappaV$.

Compare equal workloads: the same supplied source, the same necessary
preparation, the same certificate acceptance contract, and separately
reported fresh replay. Relabel every source several times to measure
early-return ordering effects. Record failures and cap exhaustion beside
successes. In particular, the elementary all-subspace method's slower
small-source performance should remain visible when evaluating the
optimized replacement.

Three falsifiable milestones would make the next iteration decisive:
\begin{enumerate}
\item General-$k$ row-first enumeration exactly matches elementary
enumeration on a complete small corpus while reducing actual attempted
systems and elapsed time.
\item A certified move family has a proved inequality for cancellation
support or gauge dimension and preserves source size within an explicit
bound. A single verified counterfamily should trigger revision of the
proposed structural claim.
\item Partial disc-base replacements output independently checked
boundary maps and an admissible next handle structure on both the
full-circle and collar fixtures, before being used inside a hierarchy.
\end{enumerate}

\section{Conclusion}

Sparse edge cancellation gives a complete way to discover standard vertex
meridians without first supplying their quadrilateral sector. Its search
spaces have dimension at most $k+1$ in the one-vertex rank-one case, and
sparse equation support gives the sharper explicit enumeration bound.
Exact rank, Euler characteristic and boundary parity suffice to verify
every promised vertex disc in polynomial bit time. The new disc and its
infinite layered extensions demonstrate that this parameter can remain
constant in the presence of dense quadrilateral support, exponential
normal weights and persistent failure of coherent-only discovery.

The Whitehead construction shows why source selection is an essential
part of the proposed global theorem. Ordered parallelity interfaces
provide an additional geometric operation needed for controlled changes
of source, but the admissible replacement and coverage arguments still
require proofs. The precise remaining objective is therefore a
constructive, certified source family with logarithmic $p+k$. Proving it,
or finding the right replacement parameter when it fails, is the central
research task exposed by this continuation.
